\daytitle{10}{冲量、动量与动量守恒}

\begin{learningbox}{本节导读}
前两节已经用“功和能量”把初末状态联系起来，但碰撞、爆炸、反冲这类过程往往时间极短，力很大且不容易直接求功。此时更适合看力在时间上的积累，也就是冲量和动量。动量守恒是继牛顿定律和能量方法之后的第三个核心工具，重点处理\textbf{碰撞、爆炸、反冲}等瞬间相互作用问题。
\end{learningbox}

\textbf{本节读题路线：} 先圈定系统，再看研究的是一个物体的动量变化，还是多个物体的总动量；若题目是碰撞、爆炸、反冲，重点检查外力冲量能否忽略。动量守恒只能给速度关系，若还要判断能量损失或弹性碰撞，还要再补上能量关系。

\section*{一、学习目标}

\begin{enumerate}
  \item 能用冲量 $I = Ft$ 和动量 $p = mv$ 描述力的时间积累与运动状态。
  \item 能由合外力冲量与动量变化推出动量定理 $F_{\text{合}}t = \Delta p$，并解释缓冲现象。
  \item 能判断动量守恒条件：系统合外力的冲量为零。
  \item 能区分弹性碰撞、完全非弹性碰撞的条件、能量特点和常用关系。
  \item 会先圈定系统并检查外力冲量，再用动量守恒解决碰撞、爆炸、反冲问题。
  \item 能说明滑块-木板模型中“动量守恒、能量损失、相对位移”的联系。
\end{enumerate}
\section*{二、初高中衔接}

初中物理对碰撞只有定性的了解——“碰撞时力很大、时间很短”。高中进一步用冲量和动量描述这类短时相互作用：

碰撞中发生了什么？
\begin{enumerate}
  \item 碰撞时间通常很短，但具体数量级取决于物体尺度、材料和碰撞速度。
  \item 碰撞力往往很大，而且通常随时间变化，不能只靠某一时刻的力描述整个过程。
  \item 若外力冲量远小于碰撞内力冲量，重力等外力的冲量通常可近似忽略。
\end{enumerate}

结论：\textbf{在外力冲量可忽略的碰撞过程中，系统的总动量近似守恒。}

\textbf{科学发展视角：} 汽车碰撞测试、运动护具设计、火箭分离、航天器对接和粒子碰撞实验，都离不开动量与冲量。工程上常通过延长碰撞时间来减小平均冲击力；在受力随时间变化规律相近时，峰值力通常也会随之降低。动量守恒还可用于反推出碰撞前后的速度和能量变化。

\section*{三、核心概念}

\begin{conceptbox}{冲量与动量}

\textbf{动量：}
\[
\boxed{p = mv}
\]
\begin{itemize}
  \item 矢量，方向与速度方向相同。
  \item 单位：$\mathrm{kg \cdot m/s}$。
\end{itemize}

\textbf{冲量：}
\[
\boxed{I = F t}
\]
\begin{itemize}
  \item 矢量，方向与力 $F$ 方向相同。
  \item 单位：$\mathrm{N \cdot s}$（与动量单位等价）。
  \item 冲量是力在时间上的积累效应；$I=Ft$ 严格适用于恒力。若力变化，$F$ 应理解为该时间内的等效平均力。
\end{itemize}

\end{conceptbox}

\begin{conceptbox}{动量定理}

\[
\boxed{I_{\text{合}}=\Delta p = p_2 - p_1 = mv_2 - mv_1}
\]
若合外力恒定，或用平均合外力表示，则可写成
\[
F_{\text{合}}t=\Delta p.
\]

\textbf{来源（牛顿第二定律的时间积累）：}
由 $F_{\text{合}}=ma=m\dfrac{\Delta v}{\Delta t}$，两边同乘 $\Delta t$，得
\[
F_{\text{合}}\Delta t = m v_2 - m v_1 = \Delta p
\]
若力随时间变化，$F_{\text{合}}t$ 应理解为合外力的总冲量（平均合力乘作用时间）。

\begin{itemize}
  \item 合外力的冲量等于物体动量的变化。
  \item 动量定理不关心中间过程，只关心初末动量。
  \item 与动能定理类似，但动量定理处理的是\textbf{时间}积累，动能定理处理的是\textbf{位移}积累。
\end{itemize}

\textbf{生活应用与工程设计：} 跳远时沙坑、汽车安全气囊、头盔、缓冲垫，都利用同一个思路：在动量变化 $\Delta p$ 基本相同的情况下，通过延长作用时间 $t$ 来减小平均作用力 $F$。

\begin{center}
\begin{minipage}{0.86\textwidth}
\centering
\begin{tikzpicture}[scale=0.72,>=stealth]
  \draw[->] (0,0) -- (6.4,0) node[right] {$t$};
  \draw[->] (0,0) -- (0,3.8) node[above left=1pt] {$F_{\text{平均}}$};

  \fill[red!25] (0.9,0) rectangle (1.65,3.15);
  \draw[thick,red] (0.9,0) rectangle (1.65,3.15);
  \node[red!80!black,align=center] at (1.28,3.45) {硬碰撞\\时间短、平均力大};

  \fill[blue!20] (3.45,0) rectangle (5.55,1.12);
  \draw[thick,blue] (3.45,0) rectangle (5.55,1.12);
  \node[blue!80!black,align=center] at (4.5,1.55) {缓冲\\时间长、平均力小};

  \node at (3.2,-0.55) {两块面积相同：$I=F_{\text{平均}}\Delta t=\Delta p$};
\end{tikzpicture}


\figcaption{冲量相同时，延长作用时间可以减小平均作用力}\label{fig:impulse-buffer}
\end{minipage}
\end{center}

\end{conceptbox}

\begin{conceptbox}{动量守恒定律}

\textbf{条件：} 系统合外力的冲量为零。系统不受外力，或合外力为零，是它的常见特例。

\textbf{定律：}
\[
\boxed{m_1 v_1 + m_2 v_2 = m_1 v_1' + m_2 v_2'}
\]

\textbf{来源（系统内力成对抵消）：}
内力成对出现，冲量等大反向，使系统内部动量变化相互抵消；若系统合外力的冲量为零，总动量就保持不变。

\begin{itemize}
  \item 系统内力（碰撞中两物体的相互作用力）不改变总动量。
  \item 动量守恒是矢量守恒——在一维问题中要规定正方向。
  \item 碰撞、爆炸、反冲的时间极短；若这段时间内系统合外力冲量远小于系统内部相互作用的冲量，可近似认为动量守恒。
\end{itemize}

\begin{center}
% TikZ 图：碰撞模型
% Day 10：冲量、动量与动量守恒
\begin{tikzpicture}[scale=0.9,>=Latex]
  % === 碰撞前 ===
  % m1
  \draw[thick] (-3,0.5) rectangle (-1.5,1.7);
  \node at (-2.25,1.1) {$m_1$};
  \draw[->,thick,red] (-1.5,1.1) -- (-0.3,1.1) node[above] {$v_1$};
  % m2
  \draw[thick] (0.5,0.5) rectangle (2.0,1.7);
  \node at (1.25,1.1) {$m_2$};
  \draw[->,thick,red] (2.5,1.1) -- (3.5,1.1) node[above] {$v_2$};
  \node[below] at (-0.75,0.3) {碰撞前};
  
  % === 碰撞后 ===
  % m1'
  \draw[thick] (-2,-2.0) rectangle (-0.5,-0.8);
  \node at (-1.25,-1.4) {$m_1$};
  \draw[->,thick,blue] (-0.5,-1.4) -- (0.8,-1.4) node[above] {$v_1'$};
  % m2'
  \draw[thick] (1.5,-2.0) rectangle (3.0,-0.8);
  \node at (2.25,-1.4) {$m_2$};
  \draw[->,thick,blue] (3.0,-1.4) -- (4.2,-1.4) node[above] {$v_2'$};
  \node[below] at (0.75,-2.2) {碰撞后};
  
  % 动量守恒公式放在右侧空白
  \node[red] at (3.8,-2.8) {$m_1 v_1 + m_2 v_2 = m_1 v_1' + m_2 v_2'$};
\end{tikzpicture}


\figcaption{碰撞前后系统的总动量守恒}\label{fig:collision-model}
\end{center}

\end{conceptbox}

\begin{conceptbox}{碰撞的分类}

\textbf{弹性碰撞：}
\begin{itemize}
  \item 动量守恒 + 动能守恒。
  \item 系统动能不损失（理想情况）。任何碰撞都满足能量守恒，区别在于动能是否转化为内能、声能等其他形式的能量。
  \item 两个方程联立可解出 $v_1'$ 和 $v_2'$。
  \item 以一维正碰、$v_2=0$ 为例，结果为：
  \[
  v_1' = \frac{m_1 - m_2}{m_1 + m_2} v_1,\qquad
  v_2' = \frac{2m_1}{m_1 + m_2} v_1
  \]
  （若 $v_2\neq 0$，公式更复杂，用动量守恒+动能守恒联立求解即可）
\end{itemize}

\textbf{完全非弹性碰撞：}
\begin{itemize}
  \item 在同一初态、同一系统且外力冲量可忽略的一维两体碰撞模型中，动量守恒，但动能不守恒（机械能损失最大）。
  \item 碰后两物体粘在一起以共同速度运动：
  \[
  m_1 v_1 + m_2 v_2 = (m_1 + m_2) v'
  \]
\end{itemize}

\textbf{非弹性碰撞：} 介于两者之间；若碰撞过程中系统合外力冲量可忽略，则动量守恒，动能部分损失。

\textbf{动量和动能的区别：}
\begin{center}
\begin{tabular}{@{}p{0.22\textwidth}p{0.30\textwidth}p{0.36\textwidth}@{}}
\toprule
物理量 & 性质 & 碰撞中的特点 \\
\midrule
动量 $p=mv$ & 矢量，方向与速度相同 & 单个物体动量可变；系统总动量需看外力冲量 \\
动能 $E_k=\frac12 mv^2$ & 标量，非负量 & 理想弹性碰撞中守恒 \\
\bottomrule
\end{tabular}
\end{center}

\end{conceptbox}

\section*{四、公式与适用条件}

\textbf{方法使用前检查：}
\begin{enumerate}
  \item 先规定正方向，一维动量方程中所有速度都要带符号。
  \item 先选系统，再判断系统合外力冲量是否可忽略。
  \item 碰撞、爆炸等短过程常用动量守恒；前后较长的滑动过程还要结合能量。
  \item 动量守恒不等于动能守恒，是否弹性碰撞要由题意另行判断。
\end{enumerate}

\begin{learningbox}{动量守恒列式前的系统分析清单}
动量守恒题不要从公式开始，而要先把“谁和谁组成系统”说清楚。
\begin{center}
\begin{tabular}{p{0.18\textwidth}p{0.74\textwidth}}
\toprule
检查项 & 具体做法 \\
\midrule
选系统 & 圈出参与短时间相互作用的物体，如两小球、子弹和木块、人和船。系统内的相互作用冲量会成对抵消。 \\
定方向 & 一维问题先规定正方向，初、末速度全部带符号；二维问题要在相互垂直的方向上分别列动量方程。 \\
查外力冲量 & 碰撞、爆炸、反冲时间很短时，重力和支持力的冲量常可忽略；若地面摩擦、墙的作用或外拉绳作用明显，就不能直接说系统动量守恒。 \\
写初末总动量 & 用“初态总动量 = 末态总动量”列式，不要只写某一个物体的动量变化。 \\
补充条件 & 弹性碰撞再加动能守恒；完全非弹性碰撞再加碰后共速；滑块-木板还要用能量损失和相对位移关系。 \\
\bottomrule
\end{tabular}
\end{center}

\textbf{来源理解：} 动量守恒不是“碰撞题专用公式”，而是系统内相互作用的冲量成对抵消后得到的结果。真正需要检查的是系统边界外有没有不可忽略的冲量；外力冲量不能忽略时，应回到动量定理。
\end{learningbox}

\begin{center}
\begin{tabular}{p{0.37\textwidth}p{0.15\textwidth}p{0.36\textwidth}}
\toprule
公式 & 意义 & 适用条件 \\
\midrule
$p = mv$ & 动量 & 高中非相对论范围 \\
$I = Ft$ & 冲量 & 恒力 \\
$I_{\text{合}} = \Delta p$ & 动量定理 & 任意过程用合外力冲量；恒定合外力时 $F_{\text{合}}t=\Delta p$ \\
$m_1 v_1 + m_2 v_2 = m_1 v_1' + m_2 v_2'$ & 动量守恒 & 系统合外力冲量为零；一维要先定正方向 \\
$\frac12 m_1 v_1^2+\frac12 m_2 v_2^2=\frac12 m_1 v_1'^2+\frac12 m_2 v_2'^2$ & 动能守恒 & 弹性碰撞；各速度均取大小 \\
$Q=f_k s_{\text{相对}}$ & 摩擦生热 & 滑动摩擦力大小恒定 \\
\bottomrule
\end{tabular}
\end{center}

\section*{五、典型模型}

\subsection*{模型 1：子弹打木块}

子弹射入木块（完全非弹性碰撞）：
\[
m v_0 = (m + M) v'
\]
碰撞中损失的动能主要转化为系统内能，并可能伴随声能：
\[
\Delta E_{k,\text{损}} = \frac12 m v_0^2 - \frac12 (m+M) v'^2
\]

\subsection*{模型 2：爆炸/反冲}

爆炸或反冲过程中，若系统合外力冲量可忽略，则系统总动量不变；同时系统内部储存的化学能、弹性势能或其他形式的能量可能转化为动能，使系统总动能增加。若爆炸前系统总动量为 $p_0$，则
\[
p_0=m_1v_1+m_2v_2.
\]
只有系统爆炸前静止时，才有
\[
0=m_1v_1+m_2v_2,
\]
此时两个部分的动量等大、方向相反。

\begin{center}
\begin{tikzpicture}[scale=0.88,>=Stealth]
  % Initial state
  \node[font=\small\bfseries] at (-2.3,1.15) {初态};
  \draw[rounded corners=2pt,thick,fill=gray!12] (-3.15,0.15) rectangle (-1.45,0.85);
  \draw[dashed] (-2.30,0.15) -- (-2.30,0.85);
  \node at (-2.72,0.50) {$m_1$};
  \node at (-1.88,0.50) {$m_2$};
  \node[font=\small] at (-2.30,-0.15) {$p_0=0$};

  \draw[->,thick] (-0.95,0.50) -- (-0.25,0.50);

  % Final state
  \node[font=\small\bfseries] at (2.15,1.15) {末态};
  \draw[rounded corners=2pt,thick,fill=red!10] (0.65,0.15) rectangle (1.45,0.85);
  \draw[rounded corners=2pt,thick,fill=blue!10] (2.85,0.15) rectangle (3.65,0.85);
  \node at (1.05,0.50) {$m_1$};
  \node at (3.25,0.50) {$m_2$};
  \draw[->,very thick,red] (0.65,0.98) -- (-0.25,0.98)
    node[midway,above] {$v_1$};
  \draw[->,very thick,blue] (3.65,0.98) -- (4.55,0.98)
    node[midway,above] {$v_2$};
  \draw[<->,thin] (1.48,0.50) -- (2.82,0.50)
    node[midway,above,fill=white,inner sep=1pt,font=\small] {内能释放};
  \node[font=\small] at (2.15,-0.15) {$m_1v_1+m_2v_2=0$};
\end{tikzpicture}


\figcaption{初态静止的两部分系统发生爆炸或反冲}\label{fig:momentum-recoil-explosion}
\end{center}

\subsection*{模型 3：人船模型}

设人从船尾走到船头，人与船的相对位移等于船长 $L$。取人相对船行走的方向为正。人在船上行走时，人与船系统水平方向动量守恒（初动量为零）：
\[
m_{\text{人}} v_{\text{人}} + m_{\text{船}} v_{\text{船}} = 0
\]
从开始到结束这段时间内，上式始终成立；把速度关系在同一段时间内累积起来，等价于说系统质心的水平位置不变，因此有位移关系：
\[
m_{\text{人}} x_{\text{人}} + m_{\text{船}} x_{\text{船}} = 0
\]
本题中人相对船的位移为船长 $L$：
\[
x_{\text{人}} - x_{\text{船}} = L
\]

联立可解得 $x_{\text{船}} = -\dfrac{m_{\text{人}}}{m_{\text{人}}+m_{\text{船}}} L$（负号表示船向反方向移动）。
图中船的位移用同一个船上点（如船尾）的初末位置来量，因此有
$L=x_{\text{人}}+|x_{\text{船}}|$。

\begin{center}
% TikZ 图：人船模型
% Day 10：动量守恒
\begin{tikzpicture}[scale=0.86,>=Latex]
  \draw[->,thick] (4.55,1.62) -- (5.35,1.62) node[right] {正方向};
  \draw[dashed] (2.25,1.35) -- (2.25,-3.35);
  \node[right] at (2.32,1.18) {质心位置不变};

  % 初态：人站在船尾
  \node[left] at (-0.10,0.25) {初态};
  \draw[thick] (0.50,0) -- (4.00,0) -- (3.60,-0.42) -- (0.90,-0.42) -- cycle;
  \draw[thick] (0.90,0) -- (0.90,0.45);
  \draw[thick] (3.60,0) -- (3.60,0.45);
  \node at (2.25,-0.22) {$M$};
  \node[below] at (0.90,-0.45) {船尾};
  \node[below] at (3.60,-0.45) {船头};
  \filldraw[black] (0.90,0.72) circle (0.13);
  \draw[thick] (0.90,0.59) -- (0.90,0.20);
  \draw[thick] (0.90,0.40) -- (0.66,0.18);
  \draw[thick] (0.90,0.40) -- (1.14,0.18);
  \node[above] at (0.90,0.90) {$m$};

  % 末态：船向左移，人到船头
  \node[left] at (-0.10,-2.05) {末态};
  \draw[thick] (-0.20,-2.30) -- (3.30,-2.30) -- (2.90,-2.72) -- (0.20,-2.72) -- cycle;
  \draw[thick] (0.20,-2.30) -- (0.20,-1.85);
  \draw[thick] (2.90,-2.30) -- (2.90,-1.85);
  \node at (1.55,-2.52) {$M$};
  \node[left=2pt] at (0.20,-2.54) {船尾};
  \node[below] at (2.90,-2.75) {船头};
  \filldraw[black] (2.90,-1.58) circle (0.13);
  \draw[thick] (2.90,-1.71) -- (2.90,-2.10);
  \draw[thick] (2.90,-1.90) -- (2.66,-2.12);
  \draw[thick] (2.90,-1.90) -- (3.14,-2.12);
  \node[above] at (2.90,-1.40) {$m$};

  % 位移尺：A 为末态船尾，B 为初态船尾/人初位置，C 为人末位置/末态船头
  \draw[dashed] (0.20,-1.78) -- (0.20,-4.08);
  \draw[dashed] (0.90,1.18) -- (0.90,-4.08);
  \draw[dashed] (2.90,-1.35) -- (2.90,-4.08);

  \draw[<->,thick,blue] (0.20,-3.18) -- (0.90,-3.18) node[midway,above=3pt,fill=white,inner sep=1pt] {$|x_{\text{船}}|$};
  \draw[<->,thick,red] (0.90,-3.52) -- (2.90,-3.52) node[midway,above=3pt,fill=white,inner sep=1pt] {$x_{\text{人}}$};
  \draw[<->,thick] (0.20,-3.94) -- (2.90,-3.94) node[midway,below=2pt] {$L$};
\end{tikzpicture}


\figcaption{人船模型：水平方向质心位置不变}\label{fig:human-boat-model}
\end{center}

\subsection*{模型 4：滑块-木板模型}

滑块以速度 $v_0$ 滑上静止木板，地面光滑，滑块与木板之间有摩擦。把滑块和木板看成一个系统时，水平方向合外力冲量为零，因此系统水平方向动量守恒；但滑动摩擦力做功会使机械能减少，转化为内能。

\begin{center}
% TikZ 图：滑块-木板模型中的动量与相对位移
% Day 10：冲量、动量与动量守恒
\begin{tikzpicture}[scale=0.82,>=Latex,line cap=round,line join=round,
  every node/.style={font=\small}]
  % Before
  \begin{scope}[shift={(-3.6,0)}]
    \draw[thick] (-2.0,-0.70) -- (2.0,-0.70);
    \draw[fill=gray!18,thick] (-1.45,-0.45) rectangle (1.45,-0.15);
    \draw[fill=blue!12,thick] (-1.25,-0.15) rectangle (-0.55,0.35);
    \node at (-0.90,0.10) {$m$};
    \node at (0,-0.31) {$M$};
    \draw[->,red!80!black,thick] (-1.23,0.54) -- (-0.10,0.54)
      node[midway,above] {$v_0$};
    \node[align=center] at (0,-1.10) {滑块有速度，木板静止};
  \end{scope}

  % Arrow between states
  \draw[->,thick] (-0.85,0.06) -- (0.85,0.06)
    node[midway,above] {摩擦作用};

  % After
  \begin{scope}[shift={(3.6,0)}]
    \draw[thick] (-2.0,-0.70) -- (2.0,-0.70);
    \draw[fill=gray!18,thick] (-1.45,-0.45) rectangle (1.45,-0.15);
    \draw[fill=blue!12,thick] (0.08,-0.15) rectangle (0.78,0.35);
    \node at (0.43,0.10) {$m$};
    \node at (0,-0.31) {$M$};
    \draw[->,red!80!black,thick] (-0.20,0.56) -- (0.95,0.56)
      node[midway,above] {$v$};
    \draw[<->,blue!75!black] (-0.90,-0.98) -- (0.43,-0.98);
    \node[blue!75!black] at (-0.23,-1.23) {$s_{\text{相对}}$};
    \node[align=center] at (0,-1.62) {滑块和木板达到共同速度};
  \end{scope}
\end{tikzpicture}


\figcaption{滑块-木板模型：共速由动量守恒决定，摩擦生热由相对位移决定}\label{fig:slider-plank-momentum}
\end{center}

若木板足够长，最终滑块与木板达到共同速度 $v$：
\[
mv_0=(m+M)v,\qquad v=\frac{m}{m+M}v_0.
\]

系统损失的机械能变为内能：
\[
Q=\frac12mv_0^2-\frac12(m+M)v^2.
\]
另一方面，摩擦生热等于滑动摩擦力乘两物体间的相对位移：
\[
Q=f_k s_{\text{相对}}=\mu_k mg\,s_{\text{相对}}.
\]

\textbf{认知提醒：} 滑块和木板受到的摩擦力大小相等、方向相反，所以摩擦力不改变系统总动量；但它们的作用点相对滑动，机械能会减少。不要把“内力不改变总动量”误读成“内力不做功”。

\section*{六、例题精讲}

\begin{examplebox}{动量定理的应用}

\textbf{题意分析：} 质量为 $m = 0.5\ \mathrm{kg}$ 的足球以 $v_0 = 10\ \mathrm{m/s}$ 的速度飞来，运动员用头顶球后球以 $v = 8\ \mathrm{m/s}$ 的速度反向飞出，头与球的接触时间 $t = 0.05\ \mathrm{s}$。求足球受到的平均合外力；若忽略这段极短时间内的重力冲量，估算头对足球的平均作用力。

\textbf{物理模型：} 动量定理，取初速度方向为正。

\textbf{解题过程：}

初动量：$p_1 = mv_0 = 0.5 \times 10 = 5\ \mathrm{kg \cdot m/s}$
末动量：$p_2 = m(-v) = 0.5 \times (-8) = -4\ \mathrm{kg \cdot m/s}$
动量变化：$\Delta p = p_2 - p_1 = -4 - 5 = -9\ \mathrm{kg \cdot m/s}$

由动量定理（这里 $F$ 表示足球所受平均合外力）：
\[
F_{\text{合}} t = \Delta p,\qquad F_{\text{合}} \times 0.05 = -9
\]
\[
F_{\text{合}} = -180\ \mathrm{N}
\]

负号表示平均合外力方向与初速度方向相反。由于接触时间很短，重力冲量相对于头的冲量较小；若忽略重力冲量，则头对足球的平均作用力约为 $180\ \mathrm{N}$，方向与初速度方向相反。

\textbf{方法总结：} 用动量定理前先规定正方向，再把初末动量写成带正负号的代数量；还要分清由 $\Delta p/\Delta t$ 直接得到的是平均合外力，求某一个力时需再结合受力分析。

\end{examplebox}

\begin{examplebox}{弹性碰撞}

\textbf{题意分析：} 质量为 $m_1 = 2\ \mathrm{kg}$ 的小球以 $v_1 = 3\ \mathrm{m/s}$ 的速度碰撞静止的 $m_2 = 1\ \mathrm{kg}$ 小球，发生弹性碰撞。求碰后两球的速度。

\textbf{物理模型：} 弹性碰撞，动量守恒 + 动能守恒。

\textbf{解题过程：}

动量守恒：
\[
2 \times 3 + 0 = 2 v_1' + 1 \times v_2'
\]
\[
6 = 2 v_1' + v_2'
\]

动能守恒：
\[
\frac12 \times 2 \times 3^2 = \frac12 \times 2 \times v_1'^2 + \frac12 \times 1 \times v_2'^2
\]
\[
9 = v_1'^2 + 0.5 v_2'^2
\]

解方程组（或用公式 $v_1' = \frac{m_1-m_2}{m_1+m_2}v_1$）：
\[
v_1' = \frac{2-1}{2+1} \times 3 = 1\ \mathrm{m/s},\qquad
v_2' = \frac{2 \times 2}{2+1} \times 3 = 4\ \mathrm{m/s}
\]

\textbf{方法总结：} 弹性碰撞公式使用前先确认“动量守恒 + 动能守恒”。算完后用两个极端情况检查：$m_1=m_2$ 时速度交换，$m_1\gg m_2$ 时 $m_1$ 几乎不变。

\end{examplebox}

\begin{examplebox}{完全非弹性碰撞——子弹打木块}

\textbf{题意分析：} 质量为 $m = 10\ \mathrm{g}$ 的子弹以 $v_0 = 400\ \mathrm{m/s}$ 射入质量为 $M = 0.99\ \mathrm{kg}$、静置在粗糙水平面上的木块并留在其中。已知木块与地面间的动摩擦因数 $\mu = 0.2$，$g = 10\ \mathrm{m/s^2}$。求：
\begin{enumerate}
  \item 子弹和木块的共同速度。
  \item 系统损失的动能。
  \item 碰后整体在水平面上滑行的距离。
\end{enumerate}

\textbf{物理模型：} 碰撞阶段时间极短，地面摩擦力的冲量远小于子弹与木块间的相互作用冲量，故水平方向近似动量守恒；碰后滑动阶段再用动能定理。

\begin{center}
\begin{tikzpicture}[scale=0.82,>=Stealth]
  % Stage 1
  \node[font=\small\bfseries] at (-2.4,1.40) {碰撞阶段（时间极短）};
  \draw[thick] (-4.20,0) -- (-0.55,0);
  \draw[fill=gray!12,thick] (-1.75,0.05) rectangle (-0.75,0.80);
  \node at (-1.25,0.43) {$M$};
  \filldraw[black] (-3.75,0.43) circle (0.07);
  \draw[->,very thick,red] (-3.62,0.43) -- (-2.35,0.43)
    node[midway,above] {$v_0$};
  \node[below,font=\small] at (-2.35,-0.08) {忽略地面摩擦冲量};

  \draw[->,thick] (-0.20,0.43) -- (0.50,0.43);

  % Stage 2
  \node[font=\small\bfseries] at (2.70,1.40) {碰后滑动阶段};
  \draw[thick] (0.85,0) -- (4.85,0);
  \draw[fill=gray!12,thick] (1.35,0.05) rectangle (2.55,0.80);
  \filldraw[black] (1.48,0.43) circle (0.07);
  \node at (2.02,0.43) {$m+M$};
  \draw[->,very thick,blue] (2.62,0.43) -- (3.68,0.43)
    node[midway,above] {$v'$};
  \draw[->,thick,red] (1.55,-0.22) -- (0.90,-0.22)
    node[midway,below] {$f_k$};
  \draw[<->,thin] (1.35,-0.72) -- (4.40,-0.72)
    node[midway,below] {$x$};
\end{tikzpicture}


\figcaption{子弹打木块要分为“短时碰撞”和“碰后滑动”两个阶段}\label{fig:bullet-block-stages}
\end{center}

\textbf{解题过程：}

\begin{enumerate}
  \item \textbf{碰撞阶段：} 忽略地面摩擦力在极短碰撞时间内的冲量，子弹与木块系统水平方向动量近似守恒：
  \[
  m v_0 = (m + M) v',\qquad 0.01 \times 400 = 1 \times v'
  \]
  \[
  v' = 4\ \mathrm{m/s}
  \]

  \item \textbf{碰撞中的动能损失：}
  \[
  \Delta E_k = \frac12 m v_0^2 - \frac12 (m+M) v'^2
  \]
  \[
  = \frac12 \times 0.01 \times 400^2 - \frac12 \times 1 \times 4^2
  \]
  \[
  = 800 - 8 = 792\ \mathrm{J}
  \]

  \item \textbf{碰后滑动阶段：} 子弹与木块已经成为整体，应用动能定理：
  \[
  -\mu (m+M) g x = 0 - \frac12 (m+M) v'^2
  \]
  \[
  x = \frac{v'^2}{2\mu g} = \frac{16}{2 \times 0.2 \times 10} = 4\ \mathrm{m}
  \]
\end{enumerate}

\textbf{方法总结：}
\begin{enumerate}
  \item 子弹打木块是典型的完全非弹性碰撞；动量守恒只用于外力冲量可忽略的短时碰撞阶段。
  \item 碰撞中损失的动能主要转化为系统内能，并可能伴随声能，不能简单等同于“摩擦生热”。
  \item 碰后较长时间的滑动过程受地面摩擦力作用，应改用动能定理或牛顿定律。
\end{enumerate}

\end{examplebox}

\begin{examplebox}{滑块-木板——动量与能量结合}

\textbf{题意分析：} 质量 $m=1\ \mathrm{kg}$ 的滑块以 $v_0=6\ \mathrm{m/s}$ 的速度滑上静止在光滑水平面上的长木板，木板质量 $M=2\ \mathrm{kg}$。滑块与木板间动摩擦因数 $\mu_k=0.3$，$g=10\ \mathrm{m/s^2}$，木板足够长。求最终共同速度和滑块相对木板滑行的距离。

\textbf{物理模型：} 滑块-木板模型。水平方向系统动量守恒；机械能损失转化为摩擦生热。

\textbf{解题过程：}

最终共速为 $v$。动量守恒：
\[
mv_0=(m+M)v,\qquad 1\times6=(1+2)v,
\]
\[
v=2\ \mathrm{m/s}.
\]

损失的机械能：
\[
Q=\frac12mv_0^2-\frac12(m+M)v^2
=\frac12\times1\times6^2-\frac12\times3\times2^2
=18-6=12\ \mathrm{J}.
\]

摩擦力大小 $f_k=\mu_kmg=0.3\times1\times10=3\ \mathrm{N}$，所以
\[
Q=f_k s_{\text{相对}},\qquad s_{\text{相对}}=\frac{12}{3}=4\ \mathrm{m}.
\]

\textbf{方法总结：} 先用动量守恒求共速，再用能量损失求相对滑动距离。若木板不够长，滑块可能先滑出木板，末态就不是共速，不能套这组结论。
\end{examplebox}

\section*{七、变式训练}

\begin{enumerate}
  \item 质量为 $m = 2\ \mathrm{kg}$ 的小滑块以 $v_0 = 6\ \mathrm{m/s}$ 的速度滑上静止在光滑水平面上的小车，小车质量为 $M = 4\ \mathrm{kg}$。滑块与小车间的动摩擦因数为 $\mu_k = 0.3$，$g=10\ \mathrm{m/s^2}$，且小车足够长。求二者的共同速度和滑块相对小车的滑行距离。
  \begin{center}
  \begin{tikzpicture}[scale=0.82,>=Stealth]
  \draw[thick] (-0.2,0) -- (5.6,0);
  \draw[fill=gray!10] (2.6,0.12) rectangle (4.45,0.75);
  \draw (2.95,0.12) circle (0.12);
  \draw (4.1,0.12) circle (0.12);
  \draw[fill=black] (0.9,0.32) circle (0.08);
  \draw (0.9,0.24) -- (0.9,0.75);
  \draw (0.7,0.52) -- (1.1,0.52);
  \draw (0.9,0.24) -- (0.7,0.02);
  \draw (0.9,0.24) -- (1.1,0.02);
  \draw[->,blue,thick] (1.12,0.62) -- (2.05,0.62) node[midway,above] {$v_0$};
  \draw[->,red,thick] (1.25,0.15) -- (0.55,0.15) node[midway,below] {$f$};
  \draw[->,red,thick] (2.85,0.5) -- (3.55,0.5);
  \node at (3.55,0.47) {$M$};
  \node at (0.45,0.75) {$m$};
  \node[font=\small] at (2.8,1.72) {人跳上车，最终共速};
\end{tikzpicture}

  \end{center}

  \item $m_1 = 4\ \mathrm{kg}$ 的小球以 $v_1 = 2\ \mathrm{m/s}$ 正面碰撞静止的 $m_2 = 2\ \mathrm{kg}$ 小球，碰后 $m_1$ 仍沿原方向以 $v_1' = 1\ \mathrm{m/s}$ 运动。求 $v_2'$，并通过比较碰撞前后的总动能判断碰撞类型。

  \item 质量为 $M = 150\ \mathrm{kg}$ 的船上站着一个质量为 $m = 60\ \mathrm{kg}$ 的人，船长 $L = 5\ \mathrm{m}$，船静止在水面上（阻力不计）。人从船尾走到船头，求船相对于岸移动的距离。
  \begin{center}
  % TikZ 图：人船模型
% Day 10：动量守恒
\begin{tikzpicture}[scale=0.86,>=Latex]
  \draw[->,thick] (4.55,1.62) -- (5.35,1.62) node[right] {正方向};
  \draw[dashed] (2.25,1.35) -- (2.25,-3.35);
  \node[right] at (2.32,1.18) {质心位置不变};

  % 初态：人站在船尾
  \node[left] at (-0.10,0.25) {初态};
  \draw[thick] (0.50,0) -- (4.00,0) -- (3.60,-0.42) -- (0.90,-0.42) -- cycle;
  \draw[thick] (0.90,0) -- (0.90,0.45);
  \draw[thick] (3.60,0) -- (3.60,0.45);
  \node at (2.25,-0.22) {$M$};
  \node[below] at (0.90,-0.45) {船尾};
  \node[below] at (3.60,-0.45) {船头};
  \filldraw[black] (0.90,0.72) circle (0.13);
  \draw[thick] (0.90,0.59) -- (0.90,0.20);
  \draw[thick] (0.90,0.40) -- (0.66,0.18);
  \draw[thick] (0.90,0.40) -- (1.14,0.18);
  \node[above] at (0.90,0.90) {$m$};

  % 末态：船向左移，人到船头
  \node[left] at (-0.10,-2.05) {末态};
  \draw[thick] (-0.20,-2.30) -- (3.30,-2.30) -- (2.90,-2.72) -- (0.20,-2.72) -- cycle;
  \draw[thick] (0.20,-2.30) -- (0.20,-1.85);
  \draw[thick] (2.90,-2.30) -- (2.90,-1.85);
  \node at (1.55,-2.52) {$M$};
  \node[left=2pt] at (0.20,-2.54) {船尾};
  \node[below] at (2.90,-2.75) {船头};
  \filldraw[black] (2.90,-1.58) circle (0.13);
  \draw[thick] (2.90,-1.71) -- (2.90,-2.10);
  \draw[thick] (2.90,-1.90) -- (2.66,-2.12);
  \draw[thick] (2.90,-1.90) -- (3.14,-2.12);
  \node[above] at (2.90,-1.40) {$m$};

  % 位移尺：A 为末态船尾，B 为初态船尾/人初位置，C 为人末位置/末态船头
  \draw[dashed] (0.20,-1.78) -- (0.20,-4.08);
  \draw[dashed] (0.90,1.18) -- (0.90,-4.08);
  \draw[dashed] (2.90,-1.35) -- (2.90,-4.08);

  \draw[<->,thick,blue] (0.20,-3.18) -- (0.90,-3.18) node[midway,above=3pt,fill=white,inner sep=1pt] {$|x_{\text{船}}|$};
  \draw[<->,thick,red] (0.90,-3.52) -- (2.90,-3.52) node[midway,above=3pt,fill=white,inner sep=1pt] {$x_{\text{人}}$};
  \draw[<->,thick] (0.20,-3.94) -- (2.90,-3.94) node[midway,below=2pt] {$L$};
\end{tikzpicture}

  \end{center}
\end{enumerate}

\begin{learningbox}{变式训练参考结果}
\begin{enumerate}[itemsep=0.45em,topsep=0.3em,parsep=0pt]
  \item 水平方向动量守恒，故共同速度
  $v=mv_0/(m+M)=2.0\ \mathrm{m/s}$。系统机械能损失
  $24\ \mathrm J=\mu_kmg\,s_{\text{相对}}$，得
  $s_{\text{相对}}=4.0\ \mathrm m$。
  \item 由动量守恒得 $v_2'=2.0\ \mathrm{m/s}$。碰前、碰后总动能分别为
  $8\ \mathrm J$ 和 $6\ \mathrm J$，故为非弹性碰撞。
  \item 设船相对岸移动距离为 $x$，由质心位置不变得
  $Mx=m(L-x)$，故 $x=mL/(M+m)=10/7\ \mathrm m$，方向与人行走方向相反。
\end{enumerate}
\end{learningbox}

\section*{八、竞赛入门补充}

\begin{competitionbox}
\textbf{恢复系数}

恢复系数 $e$ 描述碰撞前后沿碰撞线方向\textbf{分离相对速度与接近相对速度之比}：
\[
e = \frac{v_2' - v_1'}{v_1 - v_2}
\]

\begin{itemize}[nosep]
  \item $e = 1$：弹性碰撞（动能守恒）。
  \item $e = 0$：完全非弹性碰撞（碰后共速）。
  \item $0 < e < 1$：非弹性碰撞。
  \item $e > 1$：超弹性碰撞，系统有内能等转化为动能；普通被动碰撞通常不属于这种情况。
\end{itemize}

在质量和碰撞前速度都确定时，$e$ 越小通常表示动能损失越多；但 $e$ 本身不是一个与质量、初态无关的“能量损失百分比”。

\textbf{例题：} 一个小球从 $h_1$ 高处自由落下，与地面碰撞后弹起 $h_2$ 高（$h_2 < h_1$）。求恢复系数。

\begin{center}
\begin{tikzpicture}[scale=0.78,>=Stealth]
  \draw[thick] (-2.15,0) -- (2.15,0);
  \foreach \x in {-2.0,-1.7,...,2.0}
    \draw[thin] (\x,0) -- ++(-0.18,-0.16);

  \draw[dashed] (-0.95,0.05) -- (-0.95,2.65);
  \filldraw[fill=blue!12,draw=blue!70!black,thick] (-0.95,2.35) circle (0.16);
  \draw[->,very thick,blue] (-0.68,2.25) -- (-0.68,1.25)
    node[midway,right] {$v_1=\sqrt{2gh_1}$};
  \draw[<->,thin] (-1.55,0.08) -- (-1.55,2.35)
    node[midway,left] {$h_1$};

  \draw[dashed] (0.95,0.05) -- (0.95,1.95);
  \filldraw[fill=red!12,draw=red!70!black,thick] (0.95,1.65) circle (0.16);
  \draw[->,very thick,red] (1.22,0.48) -- (1.22,1.48)
    node[midway,right] {$|v_1'|=\sqrt{2gh_2}$};
  \draw[<->,thin] (0.35,0.08) -- (0.35,1.65)
    node[midway,left] {$h_2$};

  \node[font=\small] at (-0.95,2.92) {碰前下落};
  \node[font=\small] at (0.95,2.20) {碰后上升};
\end{tikzpicture}

\end{center}

\textbf{解：} 取向下为正方向。碰地前小球速度 $v_1 = \sqrt{2gh_1}$，碰后小球速度 $v_1' = -\sqrt{2gh_2}$；地面碰前速度 $v_2=0$，碰后 $v_2'=0$。

\[
e = \frac{v_2' - v_1'}{v_1 - v_2} = \frac{0 - (-\sqrt{2gh_2})}{\sqrt{2gh_1} - 0} = \sqrt{\frac{h_2}{h_1}}
\]

\textbf{竞赛提示：} 在一维两体碰撞模型中，若恢复系数已知，并且碰撞过程中外力冲量可忽略，恢复系数方程可与动量守恒方程联立求碰后速度。
\end{competitionbox}

\section*{九、易错点}

\begin{mistakebox}
\small
\begin{enumerate}[nosep]
  \item \textbf{动量守恒中忽略方向：} 列方程前先规定正方向。速度方向与正方向相反时要加负号。

  \item \textbf{动量定理中 $F_{\text{合}}t = \Delta p$ 的 $F_{\text{合}}$ 是合力：} 先画受力图，再求合力冲量；缓冲问题中不能只看安全气囊的力。

  \item \textbf{误以为碰撞中动能一定守恒：} 先判断碰撞类型。只有理想弹性碰撞才同时用动量守恒和动能守恒。

  \item \textbf{系统选择错误：} 先圈定系统，再判断合外力冲量能否忽略。只选相互作用物体时，重力是外力；把地球纳入系统时，系统边界和方程都要改变。碰撞题通常采用前一种选法，再判断短时重力冲量是否可忽略。

  \item \textbf{混淆弹性碰撞和完全非弹性碰撞：} 先从题干确认碰后状态：只有“粘在一起、共同运动”时，才写 $v'_1=v'_2$；若碰后分开，不能把两者速度直接设为相同。随后再判断是否额外给出“弹性碰撞”，只有这时才能同时使用动能守恒。

  \item \textbf{滑块-木板模型中把位移当相对位移：} 摩擦生热用的是两接触面之间的相对滑动距离，不是滑块对地位移，也不是木板对地位移。
\end{enumerate}
\end{mistakebox}

\section*{十、今日小结}
\enlargethispage{4\baselineskip}

{\small
\begin{enumerate}[nosep]
  \item 动量 $p = mv$，冲量 $I = Ft$（恒力），动量定理 $I_{\text{合}} = \Delta p$。
  \item 动量守恒条件：系统合外力冲量为零（碰撞中常可近似）。
  \item 弹性碰撞：动量守恒 $+$ 动能守恒。
  \item 完全非弹性碰撞：在外力冲量可忽略的一维两体模型中，动量守恒，碰后共速，机械能损失最大。
  \item 反冲/爆炸：系统总动量保持不变；只有初态静止时，各部分末动量的矢量和才为零。
  \item 滑块-木板模型：动量守恒求共速，机械能损失等于摩擦生热 $Q=f_k s_{\text{相对}}$。
\end{enumerate}

\textbf{一分钟复盘：} 碰撞、爆炸或反冲题先圈定系统，再判断外力冲量能否忽略、碰后是否共速、机械能是否守恒，据此选择动量与能量关系。
}
